% Justificación para objetivos.tex: redacción continua, sin subtítulos de
% impacto, con introducción breve y cierre sobre la viabilidad.
% Parte del contenido de secciones/10-justificacion-entregable.tex: si cambia
% un impacto, revisar ambos archivos.
\section{Justificación}
Un sistema de acceso con códigos QR dinámicos no solo cambia la forma en que
se entra al plantel; también tiene efectos en las personas que lo usan, en la
tecnología que se adopta, en los costos de operación y en el manejo de los
datos personales. Por ello, a continuación se describe el impacto del
proyecto en las dimensiones social, tecnológica, económica y ética.

En el ámbito social, la investigación puede beneficiar, en primer lugar, a
los estudiantes, que pierden tiempo en las filas de entrada y se exponen a
llegar tarde a su primera clase. Medir el tiempo de validación y la longitud
de la fila antes y durante el piloto permitirá saber si un cambio en el
mecanismo de acceso alivia esa situación o si la espera se debe a otras
causas. También puede beneficiar al personal de vigilancia, que revisa
cientos de credenciales en pocos minutos y cuya revisión depende solo de su
atención en los momentos de mayor afluencia. Por último, el Departamento de
Control Escolar y los directivos contarían con datos de ingreso por franja
horaria que hoy no existen, útiles para planear la apertura de accesos y el
personal necesario en cada horario. Como información complementaria, el
registro de intentos rechazados aportaría una primera medida de la
presentación de credenciales ajenas o vencidas, una preocupación de
seguridad cuya magnitud tampoco se conoce.

En cuanto al impacto tecnológico, el proyecto integra componentes de
hardware y software en un mismo sistema: un microcontrolador con cámara que
captura y decodifica el código, un controlador que libera el torniquete, un
servidor que valida el código y una base de datos con un panel web para
consultar los registros. La evidencia disponible sobre el ESP32-CAM aborda la
transmisión de video y no la decodificación de códigos QR
\parencite{nowroz2025esp32cam}, por lo que las pruebas del prototipo bajo
distintas condiciones de iluminación, brillo de pantalla y distancia
permitirán conocer su tiempo de lectura y su tasa de lecturas fallidas.
Estos resultados, junto con la documentación del diseño, podrán servir como
precedente para otros proyectos del \tescha{} que busquen automatizar
procesos con componentes de bajo costo.

Desde el punto de vista económico, el sistema se apoya en el dispositivo móvil que ya
posee el estudiante, por lo que no es necesario emitir tarjetas ni
credenciales adicionales, y el lector se construye con componentes de bajo
costo. Frente a otras tecnologías de validación, como NFC o el
reconocimiento facial, el código QR requiere menos infraestructura. Sin
embargo, cualquier ahorro en tiempo de trabajo del personal o en costos de
operación tendría que comprobarse en evaluaciones posteriores; esta
investigación solo aportará la información del piloto para valorarlo.

Respecto al impacto ético, el sistema trata datos personales de los
estudiantes, como su identidad y la hora de ingreso. Por ello, antes del
piloto se identificarán los requisitos de la legislación estatal
\parencite{edomex2017datos} y general \parencite{lgpdppso2025}, y se
definirán con Control Escolar el aviso de privacidad y el resguardo de los
registros. A diferencia del reconocimiento facial, el código QR no trata
datos biométricos, y quien no pueda mostrar su código seguirá el
procedimiento manual actual.

Además de estos impactos, el proyecto es viable de realizar. El prototipo se
construye con componentes de bajo costo y aprovecha el dispositivo móvil que los
estudiantes ya utilizan, de modo que no exige emitir nuevas credenciales ni
instalar una infraestructura compleja. A ello se suma que el piloto se
llevará a cabo en el propio acceso del \tescha{}, donde podrán medirse el
tiempo de espera y la longitud de las filas en las condiciones reales de
ingreso.

En conjunto, los resultados podrán aportar evidencia del propio plantel para
valorar si un sistema de acceso con códigos QR dinámicos ayuda a aliviar
las esperas del sistema actual y si conviene ampliar el piloto a otros
accesos. No obstante, esa posible ampliación y sus beneficios económicos
tendrían que confirmarse mediante evaluaciones posteriores.
